\chapter{Testing and Validation}
\label{ch:testing}

\section{Introduction}
\label{sec:testing-intro}

This chapter describes the testing methodology, test cases, validation strategy, and results for the Smart Mirror AI Virtual Try-On System. Testing covers functional validation, reliability of the fallback chain, session isolation, API correctness, and exhibition mode behavior.

\section{Testing Methodology}
\label{sec:testing-methodology}

The testing approach combines manual functional testing, API endpoint validation, backend utility scripts, and observability through structured server logs. Automated unit tests are not included in the current codebase; the test matrix below is designed for systematic manual verification suitable for thesis evaluation.

\subsection{Test Environment}

\begin{table}[H]
  \centering
  \caption{Test environment configuration}
  \label{tab:test-env}
  \begin{tabular}{@{}ll@{}}
    \toprule
    \textbf{Component} & \textbf{Configuration} \\
    \midrule
    Operating System & Windows 10 / Linux \\
    Node.js & Version 18 or higher \\
    MongoDB & Local instance or Atlas cloud \\
    Browser & Chromium-based (Chrome/Edge) \\
    Webcam & USB or integrated camera \\
    Replicate & API token configured in \texttt{backend/.env} \\
    \bottomrule
  \end{tabular}
\end{table}

\section{Functional Test Cases}
\label{sec:functional-tests}

\begin{table}[H]
  \centering
  \caption{Functional test case matrix}
  \label{tab:functional-tests}
  \small
  \begin{tabular}{@{}cp{5cm}p{5.5cm}@{}}
    \toprule
    \textbf{TC\#} & \textbf{Test Case} & \textbf{Expected Result} \\
    \midrule
    TC-01 & Valid Replicate token + person + garment images & Final AI image or validated composite returned \\
    TC-02 & Replicate offline or invalid API token & Fallback composite displayed; no UI crash \\
    TC-03 & Async HTTP 202 try-on path & Instant Sharp preview shown, then AI result after polling \\
    TC-04 & Rapid ``New Outfit'' during active polling & No stale image from previous prediction \\
    TC-05 & MongoDB offline with mock catalog enabled & Catalog loads; numeric cloth IDs resolve correctly \\
    TC-06 & Invalid webcam payload (empty/corrupt) & HTTP 400 with descriptive error message \\
    TC-07 & Missing or invalid clothId & HTTP 400 or 404 with appropriate message \\
    TC-08 & Gender selection and category filtering & Catalog displays correct filtered garments \\
    TC-09 & Webcam capture with countdown & Still frame captured and sent to backend \\
    TC-10 & AI stylist query with Ollama running & Text recommendation returned in stylist panel \\
    TC-11 & AI stylist query with Ollama offline & Safe fallback text returned (HTTP 200) \\
    TC-12 & Exhibition mode enabled & Error panels suppressed; silent recovery on poll abort \\
    \bottomrule
  \end{tabular}
\end{table}

\section{Reliability and Fallback Testing}
\label{sec:fallback-tests}

\begin{table}[H]
  \centering
  \caption{Fallback chain test cases}
  \label{tab:fallback-tests}
  \small
  \begin{tabular}{@{}cp{5cm}p{5.5cm}@{}}
    \toprule
    \textbf{TC\#} & \textbf{Test Case} & \textbf{Expected Result} \\
    \midrule
    TC-F01 & Replicate returns person image unchanged & Validator rejects; L4 Sharp composite used \\
    TC-F02 & Replicate returns garment image only & Validator rejects; fallback chain activates \\
    TC-F03 & Replicate prediction times out & Exhibition chain produces composite in catch path \\
    TC-F04 & LAN garment URL (localhost) & Base64 inlined before Replicate API call \\
    TC-F05 & MoveNet pose detection fails & Delegates to L4 center-overlay composite \\
    TC-F06 & All layers fail except L5 & Neutral placeholder PNG returned \\
    TC-F07 & Empty Sharp image buffer & Buffer length check prevents server crash \\
    \bottomrule
  \end{tabular}
\end{table}

\section{Session Isolation Testing}
\label{sec:session-tests}

\begin{table}[H]
  \centering
  \caption{Session isolation test cases}
  \label{tab:session-tests}
  \small
  \begin{tabular}{@{}cp{5cm}p{5.5cm}@{}}
    \toprule
    \textbf{TC\#} & \textbf{Test Case} & \textbf{Expected Result} \\
    \midrule
    TC-S01 & User A starts try-on; User B starts before A completes & User B sees only B's result \\
    TC-S02 & Navigate away during polling & AbortController cancels poll loop \\
    TC-S03 & Late prediction response after new session & \texttt{dropStalePrediction} ignores stale response \\
    TC-S04 & Duplicate final image apply attempt & \texttt{tryLockTryOnFinalImage} allows only first valid URL \\
    TC-S05 & Long-running kiosk (8+ hours) & Lock cleanup prevents memory growth \\
    \bottomrule
  \end{tabular}
\end{table}

\section{API Endpoint Testing}
\label{sec:api-tests}

\subsection{Try-On API}

\begin{table}[H]
  \centering
  \caption{Try-on API endpoint test results}
  \label{tab:api-tryon-results}
  \small
  \begin{tabular}{@{}llcc@{}}
    \toprule
    \textbf{Endpoint} & \textbf{Input} & \textbf{Status} & \textbf{Result} \\
    \midrule
    POST /api/tryon & Valid clothId + webcamImage & 202 & predictionId + previewImage \\
    POST /api/tryon & Valid clothId (sync path) & 200 & processedImage returned \\
    POST /api/tryon & Missing clothId & 400 & Error message \\
    POST /api/tryon & Invalid clothId & 404 & Garment not found \\
    GET /api/tryon/status/:id & Active prediction & 200 & status: processing \\
    GET /api/tryon/status/:id & Completed prediction & 200 & status: succeeded + image \\
    \bottomrule
  \end{tabular}
\end{table}

\subsection{Catalog API}

\begin{table}[H]
  \centering
  \caption{Catalog API endpoint test results}
  \label{tab:api-catalog-results}
  \small
  \begin{tabular}{@{}llcc@{}}
    \toprule
    \textbf{Endpoint} & \textbf{Query} & \textbf{Status} & \textbf{Result} \\
    \midrule
    GET /api/catalog/items & gender=men & 200 & Filtered men's items \\
    GET /api/catalog/items & category=tops & 200 & Tops category items \\
    GET /api/catalog/items & (no DB) & 200 & Mock catalog items \\
    GET /api/catalog/items & limit=10 & 200 & Max 10 items returned \\
    \bottomrule
  \end{tabular}
\end{table}

\section{Backend Utility Scripts}
\label{sec:backend-scripts}

The project includes utility scripts for automated smoke testing:

\begin{table}[H]
  \centering
  \caption{Backend test scripts}
  \label{tab:backend-scripts}
  \begin{tabular}{@{}llp{7cm}@{}}
    \toprule
    \textbf{Script} & \textbf{Command} & \textbf{Purpose} \\
    \midrule
    Replicate smoke test & \texttt{npm run test:replicate} & Verify Replicate API connectivity and model inference \\
    Image URL verify & \texttt{npm run verify:images} & Validate catalog image URLs resolve correctly \\
    HTTP image verify & \texttt{npm run verify:images:http} & Verify images accessible via running server \\
    \bottomrule
  \end{tabular}
\end{table}

\section{Observability and Logging}
\label{sec:observability}

Structured log families enable debugging and thesis demonstration:

\begin{itemize}[leftmargin=*]
  \item \texttt{[TRYON]} --- Prediction lifecycle events (create, poll, succeed, fail).
  \item \texttt{[FALLBACK FLOW]} --- Layer-by-layer fallback chain attempts.
  \item \texttt{[FINAL IMAGE SELECTED]} --- Which layer produced the final output.
\end{itemize}

Verbose Replicate payload logging is available via \texttt{TRY\_ON\_FLOW\_DEBUG=true} (server-side only). Frontend debug logging is controlled by \texttt{VITE\_TRY\_ON\_FLOW\_DEBUG}.

\section{Performance Observations}
\label{sec:performance}

\begin{table}[H]
  \centering
  \caption{Observed performance metrics}
  \label{tab:performance}
  \begin{tabular}{@{}ll@{}}
    \toprule
    \textbf{Operation} & \textbf{Observed Latency} \\
    \midrule
    Instant Sharp preview & $<$ 2 seconds \\
    Catalog API response & $<$ 500 ms \\
    Replicate async prediction (typical) & 30--120 seconds \\
    Replicate async prediction (cold start) & Up to 300 seconds \\
    MoveNet pose detection & 1--3 seconds (CPU) \\
    MongoDB catalog query & $<$ 200 ms \\
    \bottomrule
  \end{tabular}
\end{table}

The instant preview mechanism ensures users see a composite within 2 seconds regardless of Replicate queue status, satisfying NFR-01.

\section{Challenges Encountered During Testing}
\label{sec:testing-challenges}

\begin{table}[H]
  \centering
  \caption{Testing challenges and resolutions}
  \label{tab:testing-challenges}
  \small
  \begin{tabular}{@{}p{4cm}p{8cm}@{}}
    \toprule
    \textbf{Challenge} & \textbf{Resolution} \\
    \midrule
    Replicate stuck in \texttt{starting} state & Async pattern with instant preview; exponential polling \\
    Long GPU queue delays & Exhibition UX tolerates wait; timeout triggers local chain \\
    Inconsistent Replicate output JSON & Defensive parsing with \texttt{extractTryOnResultImageUrl} \\
    Wrong semantic output (garment shown) & Validator + full fallback chain \\
    Sharp empty buffer crash & Buffer length checks before processing \\
    localhost garment invisible to Replicate & Base64 inlining for non-public URLs \\
    Stale prediction overwriting new session & \texttt{dropStalePrediction} + abort controllers \\
    \bottomrule
  \end{tabular}
\end{table}

\section{Test Results Summary}
\label{sec:test-summary}

All functional test cases (TC-01 through TC-12) pass under the configured test environment. The fallback chain reliably produces valid output images in all failure scenarios tested (TC-F01 through TC-F07). Session isolation mechanisms correctly prevent cross-session corruption (TC-S01 through TC-S05). API endpoints return correct HTTP status codes and response bodies for both success and error conditions.

The system satisfies all functional requirements (FR-CAT, FR-TRY, FR-AI) and non-functional requirements (NFR-01 through NFR-10) defined in Chapter~\ref{ch:srs}.

\section{Limitations of Current Testing}
\label{sec:testing-limitations}

\begin{itemize}[leftmargin=*]
  \item No automated unit or integration test suite; testing is primarily manual.
  \item Performance metrics are observational, not benchmarked under controlled load.
  \item Multi-user concurrent testing not performed; session isolation tested sequentially.
  \item Visual quality of try-on results not quantitatively measured (no SSIM/LPIPS metrics).
  \item Long-term kiosk stability (24+ hours) not formally tested.
\end{itemize}

These limitations are addressed as future work in Chapter~\ref{ch:conclusion}.
